

\documentclass{uofa-eng-assignment}
\usepackage{caption}
\usepackage{hyperref}
\usepackage{tcolorbox}
\usepackage{boxedminipage}
\hypersetup{
    colorlinks = false,
    linkbordercolor = {white},
}
\captionsetup[table]{name=Tabla, labelsep=period}
\captionsetup[figure]{name=Figura, labelsep=period}

\newcommand{\hpl}{\hyperlink}
\newcommand{\hpt}{\hypertarget}
\newcommand*{\name}{Joan Sebastian Ramirez Sanchez}
\newcommand*{\id}{2223948}
\newcommand*{\course}{Introducción a la teoría de control}
\newcommand*{\assignment}[1]{Tarea #1}
\newcommand*{\dated}{\today}
\newcommand{\tf}{\textbf}

\begin{document}

\maketitle
\section*{Vehículo Aéreo No Tripulado Cuadricóptero}

\textbf{Contexto y aplicación.} 
Los cuadricópteros son plataformas de 6 grados de libertad (6-DOF) con fuerte acoplamiento no lineal. 
Su control preciso es esencial en logística autónoma, inspección de infraestructura y movilidad aérea urbana. Las no linealidades surgen de las transformaciones de orientación mediante ángulos de Euler

\subsection*{Variables del sistema}

\begin{center}
\begin{tabular}{lll}
\hline
\textbf{Símbolo} & \textbf{Unidades} & \textbf{Descripción} \\
\hline
$\phi, \theta, \psi$ & rad & Ángulos de Euler (roll, pitch, yaw) \\
$\Omega_i$ & rad/s & Velocidad angular del rotor $i$ \\
$k_T$ & N$\cdot$s$^2$ & Constante de empuje \\
$k_D$ & N$\cdot$m$\cdot$s$^2$ & Constante de par aerodinámico \\
$I_x, I_y, I_z$ & kg$\cdot$m$^2$ & Momentos de inercia principales \\
$L$ & m & Distancia del centro al rotor \\
\hline
\end{tabular}
\end{center}

\subsection*{Ecuaciones no lineales del sistema}

\begin{tcolorbox}[colback=brown!10!white, colframe=brown!60!black, title=No linealidades principales]
\begin{align*}
    \ddot{z} &= \frac{T}{m} \cos\phi \cos\theta - g 
\quad \text{(producto de trigonom.)} \\
I_x \ddot{\phi} 
&= \dot{\theta}\dot{\psi}(I_y - I_z) + \tau_\phi
\end{align*}
\end{tcolorbox}



%%%%%%%%%%%%%%%%%%%%%%%%%%%%%%%%%%%%%%%%%%%%%%%%%%%%%%%%%%%%%%%%%%
%%%%%%%%%%%%%%%%%%%%%%%%%%%%%%%%%%%%%%%%%%%%%%%%%%%%%%%%%%%%%%%%%%
%%%%%%%%%%%%%%%%%%%%%%%%%%%%%%%%%%%%%%%%%%%%%%%%%%%%%%%%%%%%%%%%%%
%%%%%%%%%%%%%%%%%%%%%%%%%%%%%%%%%%%%%%%%%%%%%%%%%%%%%%%%%%%%%%%%%%
%%%%%%%%%%%%%%%%%%%%%%%%%%%%%%%%%%%%%%%%%%%%%%%%%%%%%%%%%%%%%%%%%%

\section{Deducción matemática del sistema}
Pendiente 


%%%%%%%%%%%%%%%%%%%%%%%%%%%%%%%%%%%%%%%%%%%%%%%%%%%%%%%%%%%%%%%%%%
%%%%%%%%%%%%%%%%%%%%%%%%%%%%%%%%%%%%%%%%%%%%%%%%%%%%%%%%%%%%%%%%%%
%%%%%%%%%%%%%%%%%%%%%%%%%%%%%%%%%%%%%%%%%%%%%%%%%%%%%%%%%%%%%%%%%%
%%%%%%%%%%%%%%%%%%%%%%%%%%%%%%%%%%%%%%%%%%%%%%%%%%%%%%%%%%%%%%%%%%


\section{Punto de equilibrio}
Para hallar el punto de equilibrio del sistema, buscamos las condiciones de vuelo estacionario (hovering). 
Esto implica que el vehículo se encuentra en una posición fija, sin velocidades ni aceleraciones lineales o angulares
por lo tanto debemos tener
\begin{itemize}
    \item Aceleraciones lineales nulas: $\ddot{z}=\ddot{\phi}=\ddot{\theta}=\ddot{\psi}=0$
    \item velocidades lineales nulas: $\dot{z}=\dot{\phi}=\dot{\theta}=\dot{\psi}=0$
\end{itemize}
Además que necesitamos que la orientación del "dron" sea tal que no este desbalanceado entonces 
necesitamos que $\phi=\theta=0$ dado que son los ángulos de roll y pitch respectivamente, 
entonces el único ángulo que puede ser diferente de cero es el yaw $\psi$ 
pero dado que no afecta a la dinámica del sistema entonces podemos asumir que $\psi=0$ 
sin perdida de generalidad, por lo tanto el punto de equilibrio del sistema 
es (\textbf{explicar que es roll y pitch en la sección 1})
por lo tanto haciendo el analizis de las ecuaciones diferenciales del sistema entonces tenemos que el punto de equilibrio es
\begin{align*}
    0= \frac{T}{m} \cos 0 \cos 0 - g \implies T=mg
\end{align*}
esto para la ecuación diferencial de la aceleración lineal en el eje z, 
y para las ecuaciones diferenciales de las aceleraciones angulares entonces tenemos que
\begin{align*}
    0 = \dot{\theta}\dot{\psi}(I_y - I_z) + \tau_\phi \implies \tau_\phi=0 \\
\end{align*}
Viendolo de manera vectorial el punto de equilibrio será 
\begin{align*}
    z_e = z_0 , \quad
    \dot{z}_e = 0, \\ 
    \phi_e = 0, \quad
    \theta_e=0, \\
    \psi_e = 0, \quad
    \dot{\phi_e}=0,\\
    \dot{\theta_e}=0, \quad
    \dot{\psi_e}=0,\\
\end{align*}
%%%%%%%%%%%%%%%%%%%%%%%%%%%%%%%%%%%%%%%%%%%%%%%%%%%%%%%%%%%%%%%%%%
%%%%%%%%%%%%%%%%%%%%%%%%%%%%%%%%%%%%%%%%%%%%%%%%%%%%%%%%%%%%%%%%%%
%%%%%%%%%%%%%%%%%%%%%%%%%%%%%%%%%%%%%%%%%%%%%%%%%%%%%%%%%%%%%%%%%%
%%%%%%%%%%%%%%%%%%%%%%%%%%%%%%%%%%%%%%%%%%%%%%%%%%%%%%%%%%%%%%%%##





\section{Linealización del sistema}
Vamos a linealizar el sistema alrededor del punto de equilibrio encontrado en la sección anterior,
para esto definiremos las siguientes variables
\begin{align*}
    \delta z = z -z_e,& \quad \delta \dot{z}= \dot{z}-0, \delta \phi = \phi - 0, \quad \delta \dot{\phi} = \dot{\phi}\\
&\delta T = T-mg, \quad \delta \tau_\phi = \tau_\phi - 0,
\end{align*}
Se considerala tambien que $\delta \theta = \theta, \delta \psi = \psi, \delta \dot{\theta} = \dot{\theta}, \delta \dot{\psi} = \dot{\psi}$, Como este modelo no tiene estas ecuaciones supondremos que estan equilibrio (cero), o cualquier caso las su atribuciones serán nulas.
\begin{align*}
    \delta \dot{z} \approx \frac{\partial}{\partial T}\Big|_{z_e} \Delta T + \frac{\partial}{\partial \phi} \Big|_{z_e} \Delta \phi + \frac{\partial }{\partial \theta} \Big|_{z_e} \Delta \theta 
\end{align*}
Dónde 
\begin{align*}
    \frac{\partial}{\partial T}=\frac{1}{m} cos \phi cos \theta \Big|_{z_e}=\frac{1}{m}\\
    \frac{\partial}{\partial \phi}= -\frac{T}{m} \sin \phi \cos \theta \Big|_{z_e} = 0 \\
    \frac{\partial}{\partial \theta} = -\frac{T}{m} \cos \phi \sin \theta \Big|_{z_e} = 0 
\end{align*}
Por lo tanto
\[
\delta \ddot{z} \approx \frac{1}{m}\delta T
\]
eso para la primera ecuación diferencial, para la segunda ecuación tenemos lo siguiente 
\begin{align*}
    \delta \ddot{\phi} \approx \frac{\partial}{\partial \dot{\theta}} \Big|_{z_e} \Delta \dot{\theta} + \frac{\partial}{\partial \dot{\psi}} \Big|_{z_e} \Delta \dot{\psi} + \frac{1}{I_x} \delta \tau_\phi 
\end{align*}
y 
\begin{align*}
    \frac{\partial}{\partial \dot{\theta}} = \frac{\dot{\psi} I_y-I_z}{I_x}\Big|_{z_e} = 0, \quad \frac{\partial}{\partial \dot{\psi}}= \frac{\dot{\theta}(I_y-I_z)}{I_x}\Big|_{z_e} = 0  
\end{align*}
Resulta \[
\delta \ddot{\phi} = \frac{1}{I_x} \delta \tau_\phi\]
ahora tomando como vector de estados $x=[\delta z, \delta \dot{z}, \delta \phi, \delta \dot{\phi}]$ y el vector de control de entrada $u=[\delta T, \delta \tau_\phi]$ entonces tenemos que

\begin{align*}
    \dot{x} = \begin{bmatrix}
        \delta \dot{z} \\
        \delta \ddot{z} \\
        \delta \dot{\phi} \\
        \delta \ddot{\phi}
    \end{bmatrix}  = \begin{bmatrix}
        0 & 1 & 0 & 0 \\
        0 & 0 & 0 & 0 \\
        0 & 0 & 0 & 1 \\
        0 & 0 & 0 & 0
    \end{bmatrix}
    \begin{bmatrix}
        \delta z \\
        \delta \dot{z} \\
        \delta \phi \\
        \delta \dot{\phi}
    \end{bmatrix} + \begin{bmatrix}
        0 & 0 \\
        \frac{1}{m} & 0 \\
        0 & 0 \\
        0 & \frac{1}{I_x}
    \end{bmatrix}
    \begin{bmatrix}
        \delta T \\
        \delta \tau_\phi
    \end{bmatrix}
\end{align*}
\section{Solución del sistema linealizado}
Dado que el modelo es $\dot{x} = Ax+Bu$, entonces tenemos que la solución será de la forma
\begin{align*}
    x = e^{At} x(0) + \int_0^t e^{A(t-s)} B  u(s) ds  
\end{align*}
con \begin{align*}
    x = \begin{bmatrix}
        \delta \dot{z} \\
        \delta \ddot{z} \\
        \delta \dot{\phi} \\
        \delta \ddot{\phi}
    \end{bmatrix}, \quad A = \begin{bmatrix}
        0 & 1 & 0 & 0 \\
        0 & 0 & 0 & 0 \\
        0 & 0 & 0 & 1 \\
        0 & 0 & 0 & 0
    \end{bmatrix}
, \quad B = \begin{bmatrix}
        0 & 0 \\
        \frac{1}{m} & 0 \\
        0 & 0 \\
        0 & \frac{1}{I_x}
    \end{bmatrix}
\end{align*}
Viendo que la matriz $A$ es un matriz nilpotente entonces la forma exponencial de la matriz es más fácil
observe que: 
\begin{align*}
    A ^2 = \begin{bmatrix}
        0& 0 & 0 & 0 \\
        0& 0 & 0 & 0 \\
        0& 0 & 0 & 0 \\
        0& 0 & 0 & 0 
    \end{bmatrix}
\end{align*}
Por lo tanto la matriz exponencial de $A$ será
\begin{align*}
    e^{At} = I+ At = \begin{bmatrix}
        1 & t & 0 & 0 \\
        0 & 1 & 0 & 0 \\
        0 & 0 & 1 & t \\
        0 & 0 & 0 & 1 \\
    \end{bmatrix}
\end{align*}
en forma matricial tenemos la solución del sistema linealizado es
\begin{align*}
    \begin{bmatrix}
    \delta z(t) \\
    \delta \dot{z}(t) \\
    \delta \phi(t) \\
    \delta \dot{\phi}(t)
    \end{bmatrix}= \begin{bmatrix}
                1 & t & 0 & 0 \\
        0 & 1 & 0 & 0 \\
        0 & 0 & 1 & t \\
        0 & 0 & 0 & 1 \\
    \end{bmatrix} \begin{bmatrix}
        \delta z(0) \\
        \delta \dot{z}(0) \\
        \delta \phi(0) \\
        \delta \dot{\phi}(0)
    \end{bmatrix} + \int_0^t \begin{bmatrix}
        1 & t & 0 & 0 \\
        0 & 1 & 0 & 0 \\
        0 & 0 & 1 & t \\
        0 & 0 & 0 & 1 \\
    \end{bmatrix} \begin{bmatrix}
        0 & 0 \\
        \frac{1}{m} & 0 \\
        0 & 0 \\
        0 & \frac{1}{I_x}
    \end{bmatrix} \begin{bmatrix}
        \delta T(s) \\
        \delta \tau_\phi(s)
    \end{bmatrix} ds
\end{align*}
Y de una manera mas desglosada la solución del sistema será
\begin{align*}
    \delta z(t) = \delta z(0)+ t \delta \dot{z}(0)+ \frac{1}{m}\int_0^t (t-s)\delta T(s) ds \\
    \delta \dot{z}(t) = \delta \dot{z}(0) + \frac{1}{m}\int_0^t (t-s)\delta T(s) ds 
\end{align*}
\[
\delta \phi(t)
=
\delta \phi(0)
+
t\,\delta \dot{\phi}(0)
+
\frac{1}{I_x}
\int_{0}^{t}
(t-\tau)\,\delta \tau_{\phi}(\tau)\, d\tau,
\]

\[
\delta \dot{\phi}(t)
=
\delta \dot{\phi}(0)
+
\frac{1}{I_x}
\int_{0}^{t}
\delta \tau_{\phi}(\tau)\, d\tau.
\]
\tf{Grafica del sistema no se como hacerla}
\section{Controlabilidad del sistema}
Ver si el sistema es controlable vamos a hacer uso del teorema de Kalman y dado que nuestra matriz $A$ es
nilpotente entonces tenemos los siguiente
 \[
 \mathcal{C} = [B,AB]
 \]
 donde 
 \[
 A= \begin{bmatrix}
        0 & 1 & 0 & 0 \\
        0 & 0 & 0 & 0 \\
        0 & 0 & 0 & 1 \\
        0 & 0 & 0 & 0
    \end{bmatrix}, \quad AB = \begin{bmatrix}
        0 & 0 \\
        \frac{1}{m} & 0 \\
        0 & 0 \\
        0 & \frac{1}{I_x}
    \end{bmatrix}
 \]
 Asi $\mathcal{C}$ es 
  \[
\begin{bmatrix}
0 & 0 & \frac{1}{m} & 0 & 0 & 0 & 0 & 0 \\
\frac{1}{m} & 0 & 0 & 0 & 0 & 0 & 0 & 0 \\
0 & 0 & 0 & \frac{1}{I_x} & 0 & 0 & 0 & 0 \\
0 & \frac{1}{I_x} & 0 & 0 & 0 & 0 & 0 & 0
\end{bmatrix}
\]

 Por lo tanto $Rank(\mathcal{C})= 4$  dado que todas sus filas son linealmente independientes entonces el sistema es controlable
\section{Conjunto de datos admisibles}
Dado que no hay restricciones en las entradas, cualquier vector de estado inicial puede ser llevado al origen (o a cualquier estado 
deseado) en un tiempo finito, por lo tanto el conjunto de datos iniciales admisibles es $\mathbb{R}^4$.

\section{Control en un tiempo dado}

% \begin{align*}
%     \begin{cases}
%         x_1=z, \quad x_2=\dot{z} \\
%         x_3=\phi, \quad x_4=\dot{\phi} \\
%         x_5=\theta, \quad x_6=\dot{\theta} \\
%         x_7=\psi, \quad x_8=\dot{\psi}
%     \end{cases}    
% \end{align*}
% Por lo tanto nos queda el siguiente sistema de ecuaciones diferenciales de primer orden
% \begin{align*}
%     \begin{cases}
%         \dot{x}_1 = x_2 \\
%         \dot{x}_2 = \frac{\tau_z}{m} \cos x_3 \cos x_5 - g \\
%         \dot{x}_3 = x_4 \\
%         \dot{x}_4 = \frac{1}{I_x} \left( x_6 x_8 (I_y - I_z) + \tau_\phi \right) \\
%         \dot{x}_5 = x_6 \\
%         \dot{x}_6 = \frac{1}{I_y} \left( x_4 x_8 (I_z - I_x) + \tau_\theta \right) \\
%         \dot{x}_7 = x_8 \\
%         \dot{x}_8 = \frac{1}{I_z} \left( x_4 x_6 (I_x - I_y) + \tau_\psi \right)
%     \end{cases}
% \end{align*}
% dónde el control de entrada del sistema es el vector $\tau=[\tau_z, \tau_\phi, \tau_\theta, \tau_\psi]$ \\
% Asi vamos a escribir el sistema como $\dot{x}=f(x,u)$, 
% entonces tenemos que el sistema linealizado alrededor del punto de equilibrio 
% $x_e=[0,0,0,0,0,0,0,0]$ y $u_e=[mg,0,0,0]$ es
% \begin{align*}
%     \triangle \dot{x} \approx A \triangle x + B \triangle u
% \end{align*}
% dónde 
% \[
%  A = \frac{\partial f}{\partial x} = |_{(x_e,u_e)}, \quad B = \frac{\partial f}{\partial u} = |_{(x_e,u_e)}
% \]

% \begin{align*}
%     \frac{\partial f_1}{\partial x_1} = 0, \quad \frac{\partial f_1}{\partial x_2} = 1, \quad \frac{\partial f_1}{\partial x_i} = 0 \text{ para } i=3,...,8 
% \end{align*}
% para $f_2$ 
% \begin{align*}
%      \frac{\partial f_2}{\partial x_3} = -\frac  {mg}{m} \sin 0 \cos 0 = 0, \quad \frac{\partial f_2}{\partial x_5} = -\frac{mg}{m} \cos 0 \sin 0 = 0, \quad \frac{\partial f_2}{\partial x_i} = 0 \text{ para } i=1,2,4,6,7,8 \\
% \end{align*}
% para $f_3$
% \begin{align*}
%     \frac{\partial f_3}{\partial x_3} = 0, \quad \frac{\partial f_3}{\partial x_4} = 1, \quad \frac{\partial f_3}{\partial x_i} = 0 \text{ para } i=1,2,5,6,7,8
% \end{align*}
% para $f_4$
% \begin{align*}
%     \frac{\partial f_4}{\partial x_4} = 0, \quad \frac{\partial f_4}{\partial x_6} = \frac{1}{I_x} (I_y - I_z), \quad \frac{\partial f_4}{\partial x_8} = 0, \quad \frac{\partial f_4}{\partial x_i} = 0 \text{ para } i=1,2,3,5,7
% \end{align*}
% para $f_5$
% \begin{align*}
%     \frac{\partial f_5}{\partial x_5} = 0, \quad \frac{\partial f_5}{\partial x_6} = 1, \quad \frac{\partial f_5}{\partial x_i} = 0 \text{ para } i=1,2,3,4,7,8 
% \end{align*}
% para $f_6$
% \begin{align*}
%     \frac{\partial f_6}{\partial x_4} = 0, \quad \frac{\partial f_6}{\partial x_6} = 0, \quad \frac{\partial f_6}{\partial x_8} = \frac{1}{I_y} (I_z - I_x), \quad \frac{\partial f_6}{\partial x_i} = 0 \text{ para } i=1,2,3,5,7
% \end{align*}
% para $f_7$
% \begin{align*}
%     \frac{\partial f_7}{\partial x_7} = 0, \quad \frac{\partial f_7}{\partial x_8} = 1, \quad \frac{\partial f_7}{\partial x_i} = 0 \text{ para } i=1,2,3,4,5,6
% \end{align*}
% para $f_8$
% \begin{align*}
%     \frac{\partial f_8}{\partial x_4} = 0, \quad \frac{\partial f_8}{\partial x_6} = 0, \quad \frac{\partial f_8}{\partial x_8} = 0, \quad \frac{\partial f_8}{\partial x_i} = 0 \text{ para } i=1,2,3,5,7
% \end{align*}
% entonces la matriz $A$ es
% \[
% A = \begin{bmatrix}
% 0 & 1 & 0 & 0 & 0 & 0 & 0 & 0 \\
% 0 & 0 & 0 & 0 & 0 & 0 & 0 & 0 \\
% 0 & 0 & 0 & 1 & 0 & 0 & 0 & 0 \\
% 0 & 0 & 0 & 0 & 0 & 0& 0 & 0  \\
% 0 & 0 & 0 & 0 & 0 & 1 & 0 & 0 \\
% 0 & 0 & 0 & 0 & 0 & 0 & 0 & 0 \\
% 0 & 0 & 0 & 0 & 0 & 0 & 0 & 1 \\
% 0 & 0 & 0 & 0 & 0 & 0 & 0 & 0 
% \end{bmatrix}
% \]
% y la matriz $B$ es
% \[
% B = \begin{bmatrix}
% 0 & 0 & 0 & 0 \\
% \frac{1}{m} & 0 & 0 & 0 \\
% 0 & 0 & 0 & 0 \\
% 0 & \frac{1}{I_x} & 0 & 0 \\
% 0 & 0 & 0 & 0 \\
% 0 & 0 & \frac{1}{I_y} & 0 \\
% 0 & 0 & 0 & 0 \\
% 0 & 0 & 0 & \frac{1}{I_z}
% \end{bmatrix}
% \]
% por lo tanto la solución del sistema linealizado es
% \begin{align*}
%     x = e^{At} C_1 + \int_0^t e^{A(t-s)} B  u(s) ds  
% \end{align*}
% Calculemos primero lo que esta dentro de la integral, donde $e^A(t-s)$ al ser una matriz nipotente
% es facil encontrar su exponencial, dado que $A^2=0$ entonces tenemos que
% \begin{align*}
%     e^{A(t-s)} = I + A(t-s)
% \end{align*}
% por lo tanto
% \[
% x = e^{At} C_1 + \int_0^t (I + A(t-s)) B u(s) ds
% \]
% \[
% x= e^{At}C_1 + \int_{0}^{t}(B+A(t-s)B) u(s)ds
% \]
% \begin{align*}
%     e^{A(t-s)} B u(s) &= \left(\begin{matrix}
% 1 & t-s & 0 & 0 & 0 & 0 & 0 & 0 \\
% 0 & 1 & 0 & 0 & 0 & 0 & 0 & 0 \\
% 0 & 0 & 1 & t-s & 0 & 0 & 0 & 0 \\
% 0 & 0 & 0 & 1 & 0 & 0 & 0 & 0 \\
% 0 & 0 & 0 & 0 & 1 & t-s & 0 & 0 \\
% 0 & 0 & 0 & 0 & 0 & 1 & 0 & 0 \\
% 0 & 0 & 0 & 0 & 0 & 0 & 1 & t-s \\
% 0 & 0 & 0 & 0 & 0 & 0 & 0 & 1
% \end{matrix}\right)
% \left(\begin{matrix}
% 0 & 0 & 0 & 0 \\
% \frac{1}{m} & 0 & 0 & 0 \\
% 0 & 0 & 0 & 0 \\
% 0 & \frac{1}{I_x} & 0 & 0 \\
% 0 & 0 & 0 & 0 \\
% 0 & 0 & \frac{1}{I_y} & 0 \\
% 0 & 0 & 0 & 0 \\
% 0 & 0 & 0 & \frac{1}{I_z}
% \end{matrix}\right) 
% \left(\begin{matrix}
% u_1(s) \\
% u_2(s) \\
% u_3(s) \\
% u_4(s)
% \end{matrix}\right)
% \\
% &= \left(\begin{matrix}
%     1 & t-s & 0 & 0 & 0 & 0 & 0 & 0 \\
% 0 & 1 & 0 & 0 & 0 & 0 & 0 & 0 \\
% 0 & 0 & 1 & t-s & 0 & 0 & 0 & 0 \\
% 0 & 0 & 0 & 1 & 0 & 0 & 0 & 0 \\
% 0 & 0 & 0 & 0 & 1 & t-s & 0 & 0 \\
% 0 & 0 & 0 & 0 & 0 & 1 & 0 & 0 \\
% 0 & 0 & 0 & 0 & 0 & 0 & 1 & t-s \\
% 0 & 0 & 0 & 0 & 0 & 0 & 0 & 1
% \end{matrix}\right)
% \left(\begin{matrix}
% 0 \\
% \frac{1}{m} u_1(s) \\
% 0 \\
% \frac{1}{I_x} u_2(s) \\
% 0 \\
% \frac{1}{I_y} u_3(s) \\ 
% 0 \\
% \frac{1}{I_z} u_4(s)
% \end{matrix}\right) \\
% &= \left(\begin{matrix}
% \frac{t-s}{m} u_1(s)  \\
% \frac{1}{m} u_1(s) \\
% \frac{t-s}{I_x} u_2(s)\\
% \frac{1}{I_x} u_2(s) \\
% \frac{t-s}{I_y} u_3(s) \\
% \frac{1}{I_y} u_3(s) \\
% \frac{t-s}{I_z} u_4(s) \\
% \frac{1}{I_z} u_4(s)
% \end{matrix}\right)
% % &= \left(\begin{matrix}
% % \frac{t-s}{m} u_1(s) + \frac{t-s}{I_x} u_2(s) + \frac{t-s}{I_y} u_3(s) + \frac{t-s}{I_z} u_4(s) \
% % \frac{1}{m} u_1(s) \\
% % \frac{t-s}{I_x} u_2(s) + \frac{t-s}{I_y} u_3(s) + \frac{t-s}{I_z} u_4(s) \\
% % \frac{1}{I_x} u_2(s) \\
% % \frac{t-s}{I_y} u_3(s) + \frac{ t-s}{I_z} u_4(s) \\
% % \frac{1}{I_y} u_3(s) \\
% % \frac{t-s}{I_z} u_4(s) \\
% % \frac{1}{I_z} u_4(s)
% % \end{matrix}\right) 
% \end{align*}
% Algo un poco feo entonces vamos a renombrar las constantes para que se vea un poco mas bonito, entonces tenemos que dónde $a_1=\frac{1}{m}, a_2=\frac{1}{I_x}, a_3=\frac{1}{I_y}, a_4=\frac{1}{I_z}$, entonces,
% \begin{align*}
%     e^{A(t-s)} B u(s) = \left(\begin{matrix}
%         a_1 (t-s) u_1(s)\\
% a_1 u_1(s) \\
% a_2 (t-s) u_2(s) \\
% a_2 u_2(s) \\
% a_3 (t-s) u_3(s) \\
% a_3 u_3(s) \\
% a_4 (t-s) u_4(s) \\
% a_4 u_4(s)
% \end{matrix}\right)
% \end{align*}
% Por lo tanto la solución del sistema linealizado es
% \begin{align*}
%     x(t) &= \left(\begin{matrix}
% 1 & t & 0 & 0 & 0 & 0 & 0 & 0 \\
% 0 & 1 & 0 & 0 & 0 & 0 & 0 & 0 \\
% 0 & 0 & 1 & t & 0 & 0 & 0 & 0 \\
% 0 & 0 & 0 & 1 & 0 & 0 & 0 & 0 \\
% 0 & 0 & 0 & 0 & 1 & t & 0 & 0 \\
% 0 & 0 & 0 & 0 & 0 & 1 & 0 & 0 \\
% 0 & 0 & 0 & 0 & 0 & 0 & 1 & t \\
% 0 & 0 & 0 & 0 & 0 & 0 & 0 & 1
% \end{matrix}\right)
%      C_1 + \int_0^t \left(\begin{matrix}
%         a_1 (t-s) u_1(s)\\
% a_1 u_1(s) \\
% a_2 (t-s) u_2(s) \\
% a_2 u_2(s) \\
% a_3 (t-s) u_3(s) \\
% a_3 u_3(s) \\
% a_4 (t-s) u_4(s) \\
% a_4 u_4(s)
% \end{matrix}\right) ds \\
% &= \left(\begin{matrix}
% C_{11} + C_{12} t + a_1 \int_0^t (t-s) u_1(s) ds \\
% C_{12} + a_1 \int_0^t u_1(s) ds \\
% C_{13} + C_{14} t + a_2 \int_0^t (t-s) u_2(s) ds \\
% C_{14} + a_2 \int_0^t u_2(s) ds \\
% C_{15} + C_{16} t + a_3 \int_0^t (t-s) u_3(s) ds \\
% C_{16} + a_3 \int_0^t u_3(s) ds \\
% C_{17} + C_{18} t + a_4 \int_0^t (t-s) u_4(s) ds \\
% C_{18} + a_4 \int_0^t u_4(s) ds
% \end{matrix}\right)
% \end{align*}
% y la solución sin el termino forzante nos queda 
% \begin{align*}
%     x(t) = \left(\begin{matrix}
%         C_{11} + C_{12} t \\
%         C_{12} \\
%         C_{13} + C_{14} t \\
%         C_{14} \\     
%         C_{15} + C_{16} t \\
%         C_{16} \\
%         C_{17} + C_{18} t \\
%         C_{18}
%     \end{matrix}\right)
% \end{align*}
% \section{Uso del teorema de Kalman}
% Vamos a ver si el sistema es controlable, para eso vamos a usar la condición de controlabilidad de Kalman, entonces tenemos que la matriz de controlabilidad es
% \begin{align*} 
%    C = Rank[B, AB, A^2 B]
% \end{align*}
% Dado que $A^2=0$ entonces $A^2 B=0$ entonces la matriz de controlabilidad es
% \begin{align*}
%     C = Rank[B, AB] = Rank\left[\begin{matrix}
% 0 & 0 & 0 & 0 \\
% \frac{1}{m} & 0 & 0 & 0 \\
% 0 & 0 & 0 & 0 \\
% 0 & \frac{1}{I_x} & 0 & 0 \\
% 0 & 0 & 0 & 0 \\  
% 0 & 0 & \frac{1}{I_y} & 0 \\
% 0 & 0 & 0 & 0 \\
% 0 & 0 & 0 & \frac{1}{I_z}
% \end{matrix}\right]
% \end{align*}
% Como $AB$ es
% \begin{align*}
%     \left(\begin{matrix}
% \frac{1}{m} & 0 & 0 & 0 \\
% 0 & 0 & 0 & 0 \\
% 0 & \frac{1}{I_x} & 0 & 0 \\
% 0 & 0 & 0 & 0 \\
% 0 & 0 & \frac{1}{I_y} & 0 \\
% 0 & 0 & 0 & 0 \\
% 0 & 0 & 0 & \frac{1}{I_z} \\
% 0 & 0 & 0 & 0
% \end{matrix}\right)
% \end{align*}
% por lo tanto la matriz de controlabilidad es
% \begin{align*}
%     C = Rank[B, AB] = 8
% \end{align*}
% Por lo tanto el sistema es controlable.
% \section{Conjunto de datos iniciales admisibles}
% Dado que al linealizar el sistema 






\end{document}
